\documentclass[11pt,a4paper]{article}
\usepackage[top=1.7cm,bottom=1.5cm,left=1.5cm,right=1.5cm,headsep=8pt,footskip=16pt]{geometry}
\usepackage{enumitem}
\usepackage{amssymb}
\usepackage{array}
\usepackage[table]{xcolor}
\usepackage{fancyhdr}
\usepackage{lastpage}
\usepackage{multicol}
\usepackage{polyglossia}
\setdefaultlanguage[calendar=gregorian, numerals=maghrib]{arabic}
\setotherlanguage{english}
\IfFontExistsTF{Amiri}{%
  \newfontfamily\arabicfont[Script=Arabic, Scale=1.08]{Amiri}
}{%
  \newfontfamily\arabicfont[Script=Arabic, Scale=1.08, Path=../fonts/, Extension=.ttf,
    UprightFont=*-Regular, BoldFont=*-Bold, ItalicFont=*-Italic,
    BoldItalicFont=*-BoldItalic]{Amiri}
}
\IfFontExistsTF{AlBattar}{%
  \newfontfamily\khatt[Script=Arabic, Scale=1.15]{AlBattar}
}{%
  \newfontfamily\khatt[Script=Arabic, Scale=1.15, Path=../fonts/,
    BoldFont=ae_AlBattar.ttf]{ae_AlBattar.ttf}
}
\renewcommand{\headrulewidth}{0.8pt}
\renewcommand{\footrulewidth}{0.8pt}
\fancyhf{}
\rhead{\khatt المدرسة العليا للأساتذة ببوزريعة}
\lhead{\khatt ملحقة بومرداس}
\lfoot{\khatt الأستاذ: د. بن عبد الرزاق يوسف}
\cfoot{\thepage/\pageref{LastPage}}
\pagestyle{fancy}
\parindent 0pt
\setlength{\columnsep}{0.6cm}
\setlength{\columnseprule}{0.4pt}
\setlist[enumerate,1]{label=\textbf{\arabic*.}, leftmargin=*, itemsep=5pt, topsep=2pt}
\begin{document}

\linespread{1.08}\selectfont
\begin{center}
{\LARGE\khatt\textbf{الاختبار السابع في الذكاء الاصطناعي}}\\[4pt]
{\normalsize وضعيات إشكالية ــ أسئلة متعددة الاختيارات ــ الفصول من 1 إلى 6}
\end{center}
\vspace{-2pt}\rule{\textwidth}{0.6pt}\\[2pt]
\textbf{التعليمة:} يتضمن الاختبار خمس وضعيات إشكالية. اقرأ كل وضعية بعناية، ثم اختر التحليل أو التصرف الأنسب
من بين الاختيارات الخمسة؛ لكل وضعية إجابة صحيحة واحدة فقط.
اكتب \textbf{حرف} الإجابة الصحيحة في جدول الإجابات في آخر الاختبار.
لكل وضعية 4 نقاط (المجموع 20 نقطة).
\vspace{4pt}

\noindent
{\large\khatt\textbf{الوضعية 1}}\hfill{\small(الفصل 1)}\par\nopagebreak\vspace{3pt}
\setlength{\fboxsep}{7pt}\colorbox{gray!10}{\parbox{\dimexpr\linewidth-14pt}{عرض مفتش على أساتذة أربع مهام: اقتراح موضوعات للتعبير الكتابي، والحكم على صحة حديث نبوي يُستشهد به في درس، وصياغة مسودة رسالة إلى الأولياء، والفصل في علامة تلميذ يطعن فيها. وطلب منهم تصنيف المهام حسب ملاءمتها للذكاء الاصطناعي التوليدي. واختلف الأساتذة في تصنيف بعضها. أي تصنيف هو الأدق؟}}\par\nopagebreak
\begin{itemize}[leftmargin=1.0cm, labelsep=6pt, itemsep=3pt, topsep=4pt, parsep=0pt]
  \item[\textbf{أ)}] اقتراح الموضوعات وحده مناسب، أما المسودة فلا، لأن الرسائل الرسمية يجب أن تُكتب يدويًا من البداية، والحكم على الحديث مناسب للأداة، لأنها اطلعت على كتب الحديث كلها.
  \item[\textbf{ب)}] الحكم على صحة الحديث والفصل في العلامة أنسب المهام للأداة لأنها محايدة ولا تتأثر بالعاطفة، أما اقتراح الموضوعات وصياغة المسودة فيتطلبان إبداعًا إنسانيًا لا تملكه الأدوات.
  \item[\textbf{ج)}] اقتراح الموضوعات وصياغة المسودة مناسبان مع المراجعة، أما صحة الحديث فيُرجع فيها إلى كتب التخريج والمختصين، والفصل في العلامة قرار تربوي يبقى للأستاذ واللجنة، ولا يُسند إلى الأداة.
  \item[\textbf{د)}] اقتراح الموضوعات وصياغة المسودة والحكم على الحديث مناسبة للأداة مع مراجعة سريعة، وحده الفصل في العلامة يبقى للأستاذ، لأنه يمس حق التلميذ، أما الحديث فمعلومة ثابتة لا تخطئ فيها الأدوات.
  \item[\textbf{هـ)}] المهام الأربع مناسبة للأداة إذا كان الطلب محكمًا، لأن جودة الطلب تضمن جودة المخرج في كل المجالات، بما فيها تخريج الأحاديث والفصل في الطعون، ما دام الأستاذ يطّلع على النتيجة.
\end{itemize}
\par\vspace{6pt}
\noindent
{\large\khatt\textbf{الوضعية 2}}\hfill{\small(الفصل 3)}\par\nopagebreak\vspace{3pt}
\setlength{\fboxsep}{7pt}\colorbox{gray!10}{\parbox{\dimexpr\linewidth-14pt}{طلبت أستاذة نصًا قصيرًا لتلاميذ السنة الأولى ابتدائي وكتبت: «لا تستعمل كلمات صعبة، ولا جملًا طويلة، ولا أسلوبًا معقدًا». فجاء النص أبسط قليلًا، لكنه بقي فوق مستوى التلاميذ في المفردات وطول الجمل. وتريد نصًا يقرؤه التلاميذ وحدهم. وجرّبت أيضًا إضافة دور «خبير» دون جدوى. أي إعادة صياغة للقيود هي الأنجع؟}}\par\nopagebreak
\begin{itemize}[leftmargin=1.0cm, labelsep=6pt, itemsep=3pt, topsep=4pt, parsep=0pt]
  \item[\textbf{أ)}] تحذف القيود كلها وتكتب «نص لطفل صغير»، لأن القيود الكثيرة تربك الأداة، والعبارة العامة القصيرة تكفي لتحديد المستوى بدقة، فالأداة تعرف ما يناسب كل الأعمار دون أي تفصيل.
  \item[\textbf{ب)}] تكرر النفي بصيغ أشد: «لا تستعمل أبدًا أي كلمة صعبة مهما كانت، ولا أي جملة طويلة إطلاقًا»، لأن تشديد النفي وتكراره يجعل الأداة أكثر التزامًا بالقيود من أي معايير تفصيلية أخرى.
  \item[\textbf{ج)}] تضيف «أنت خبير عالمي في أدب الأطفال» في بداية الطلب وتُبقي القيود المنفية كما هي، لأن تحديد الدور وحده يضبط مستوى المفردات وطول الجمل، ويغني عن أي معايير قابلة للقياس.
  \item[\textbf{د)}] تطلب النص بمستوى السنة الخامسة ابتدائي ثم تختصره بنفسها، لأن الأداة لا تستطيع الكتابة لمستوى أدنى من ذلك، والاختصار اليدوي يجعل أي نص مناسبًا لتلاميذ السنة الأولى تلقائيًا.
  \item[\textbf{هـ)}] تحوّل القيود إلى معايير محددة: «جمل لا تتجاوز ست كلمات، بأفعال مضارعة، ومفردات من محيط الطفل كالبيت والمدرسة، في ثمانية أسطر»، وتضيف فقرة من كتاب التلميذ مثالًا على المستوى المطلوب.
\end{itemize}
\par\vspace{6pt}
\noindent
{\large\khatt\textbf{الوضعية 3}}\hfill{\small(الفصل 4)}\par\nopagebreak\vspace{3pt}
\setlength{\fboxsep}{7pt}\colorbox{gray!10}{\parbox{\dimexpr\linewidth-14pt}{لاحظت أستاذة مكوِّنة أن طلبة التربص الميداني يحضّرون مذكراتهم البيداغوجية كلها بأداة ذكاء اصطناعي، فتأتي منظمة أنيقة، لكنهم يعجزون عن تعليل اختياراتهم أمام لجنة التقويم، ولا يعدّلون الخطة عند تغيّر ظروف القسم. ويرى بعض الزملاء أن الحل في المنع التام، ويرى آخرون أن الحل في الإفصاح وحده. ما المعالجة الأنسب تربويًا وأخلاقيًا؟}}\par\nopagebreak
\begin{itemize}[leftmargin=1.0cm, labelsep=6pt, itemsep=3pt, topsep=4pt, parsep=0pt]
  \item[\textbf{أ)}] لا تمنع الأداة، بل تشترط أن يُعدّ الطالب تصوره أولًا، ثم يستعين بها للتحسين، ويُفصح عن دورها، ويعلّل كل اختيار شفهيًا؛ لأن الهدف بناء كفاءة التخطيط لا إنتاج مذكرات منظمة فقط.
  \item[\textbf{ب)}] تشترط الإفصاح عن استعمال الأداة فقط، دون تغيير طريقة الإعداد أو التقويم، لأن الإفصاح يحل المشكلة الأخلاقية والتربوية معًا، ويضمن تلقائيًا أن يكتسب الطالب كفاءة التخطيط المطلوبة، دون حاجة إلى أي إجراء إضافي.
  \item[\textbf{ج)}] تسمح بالاستعمال الحر دون إفصاح، وتطلب من الطلبة حفظ تعليلات تولّدها الأداة لكل مذكرة، لأن حفظ التعليلات يحقق المطلوب أمام اللجنة، ويوفر وقت التكوين لأنشطة أخرى.
  \item[\textbf{د)}] تمنع الأداة منعًا كليًا في التكوين، لأن أي استعمال لها يُضعف كفاءة التخطيط حتمًا، ولا يمكن الجمع بين بناء الكفاءة والاستعانة بالأدوات في أي مرحلة من مراحل التكوين المهني، كما تفعل بعض مؤسسات التكوين.
  \item[\textbf{هـ)}] تكتفي بتقييم المذكرات المكتوبة وتلغي مناقشة اللجنة، لأن المذكرة المنظمة دليل كافٍ على الكفاءة، وعجز الطالب عن التعليل الشفهي مسألة تواصلية لا علاقة لها بالتخطيط البيداغوجي.
\end{itemize}
\par\vspace{6pt}
\noindent
{\large\khatt\textbf{الوضعية 4}}\hfill{\small(الفصل 5)}\par\nopagebreak\vspace{3pt}
\setlength{\fboxsep}{7pt}\colorbox{gray!10}{\parbox{\dimexpr\linewidth-14pt}{استعمل طالب أداة ذكاء اصطناعي متصلة بالبحث في الإنترنت، فأجابت عن سؤاله حول تاريخ أول مطبعة عربية مع ثلاث إحالات إلى روابط. فتحها فوجد الأول يؤكد المعلومة، والثاني يتحدث عن موضوع قريب دون أن يذكرها، والثالث معطّلًا. وموعد تسليم البحث بعد يومين. كيف يتعامل مع الإجابة في بحثه؟}}\par\nopagebreak
\begin{itemize}[leftmargin=1.0cm, labelsep=6pt, itemsep=3pt, topsep=4pt, parsep=0pt]
  \item[\textbf{أ)}] يعتمد المعلومة ويحيل إلى المصدرين الأول والثاني فقط، لأن الرابط المعطل وحده هو المشكوك فيه، أما الثاني فقريب الموضوع، وقرب الموضوع كافٍ للإحالة في البحوث الجامعية.
  \item[\textbf{ب)}] يرفض المعلومة كليًا، لأن وجود رابط معطل ورابط غير مطابق يعني أن الأداة اختلقت الإجابة كلها، حتى المعلومة التي يؤكدها المصدر الأول، فلا يجوز الاعتماد على أي جزء منها.
  \item[\textbf{ج)}] يعتمد المعلومة ويحيل إلى الأداة نفسها بدل الروابط، لأن الأداة جمعت المصادر ولخّصتها، فتكون هي المرجع الأصلي، والإحالة إليها أدق من الإحالة إلى المواقع التي استقت منها، وهذا أوضح للقارئ وأيسر في التوثيق.
  \item[\textbf{د)}] يعتمد المعلومة بإحالتها إلى المصدر الأول بعد تقييم موثوقيته، ولا ينسبها إلى الثاني لأنه لا يتضمنها، ويبحث عن مصدر علمي إضافي يؤكدها، لأن وجود رابط لا يعني أن مضمونه يدعم الإجابة.
  \item[\textbf{هـ)}] يعتمد المعلومة ويحيل إلى الروابط الثلاثة، لأن الأداة لا تضع رابطًا إلا إذا كان يدعم إجابتها، والرابط المعطل كان يتضمن المعلومة قبل تعطله، والثاني يدعمها ضمنيًا بحكم قرب موضوعه، كما هو معمول به في البحوث السريعة.
\end{itemize}
\par\vspace{6pt}
\noindent
{\large\khatt\textbf{الوضعية 5}}\hfill{\small(الفصل 6)}\par\nopagebreak\vspace{3pt}
\setlength{\fboxsep}{7pt}\colorbox{gray!10}{\parbox{\dimexpr\linewidth-14pt}{سجّلت أستاذة قراءات تلاميذها الجهرية، واستعملت أداة تحويل الصوت إلى نص لرصد أخطائهم آليًا. لاحظت أن الأداة «تصحّح» أحيانًا أخطاء التلاميذ فتكتب الكلمة الصحيحة بدل ما نطقوه فعلًا، وتخطئ مع بعض اللهجات. وتريد بناء تقويم للقراءة على هذه النتائج. وقد سبق أن قدّمت نتائج أولية لزملائها اعتمادًا على الأداة. ما الاستنتاج الأنسب؟}}\par\nopagebreak
\begin{itemize}[leftmargin=1.0cm, labelsep=6pt, itemsep=3pt, topsep=4pt, parsep=0pt]
  \item[\textbf{أ)}] المشكلة تقتصر على اللهجات، فتعتمد الأداة في تقويم التلاميذ الذين يقرؤون بنطق قريب من الفصحى، ولا تراجع تسجيلاتهم، وتستمع فقط إلى تسجيلات التلاميذ ذوي اللهجات المختلفة، فهي الحالات الوحيدة التي يقع فيها الخطأ عادة.
  \item[\textbf{ب)}] الأداة تميل إلى كتابة الكلمات المرجّحة لغويًا لا ما نُطق فعلًا، فلا تصلح وحدها لرصد أخطاء القراءة؛ يمكن استعمالها مساعدًا أوليًا، مع الاستماع إلى التسجيلات للتحقق من الأخطاء قبل أي تقويم.
  \item[\textbf{ج)}] الأدوات الصوتية لا تفهم العربية إطلاقًا، وكل ما كتبته مولَّد عشوائيًا، فتتخلى عن تسجيل القراءات أصلًا، لأن التسجيل لا فائدة منه إذا لم تستطع الأداة تحويله إلى نص صحيح.
  \item[\textbf{د)}] الأداة أدق من أذن الأستاذة، فما كتبته هو ما نطقه التلاميذ فعلًا، والأخطاء التي سمعتها الأستاذة سببها ضعف جودة التسجيل، فتعتمد النصوص المكتوبة أساسًا وحيدًا للتقويم.
  \item[\textbf{هـ)}] تطلب من التلاميذ القراءة ببطء شديد وبصوت مرتفع، لأن ذلك يمنع الأداة نهائيًا من تصحيح الكلمات تلقائيًا، فتصبح صالحة لرصد كل أخطاء القراءة دون حاجة إلى أي استماع بشري، وهي طريقة يوصي بها بعض المختصين.
\end{itemize}
\par\vspace{6pt}
\begin{center}
{\khatt\textbf{جدول الإجابات}}\\[3pt]
\renewcommand{\arraystretch}{1.6}
\begin{tabular}{|>{\centering\arraybackslash}m{1.9cm}|*{5}{>{\centering\arraybackslash}m{1.2cm}|}}
\hline
\textbf{رقم السؤال} & \textbf{1} & \textbf{2} & \textbf{3} & \textbf{4} & \textbf{5} \\ \hline
\textbf{الإجابة} &  &  &  &  &  \\[8pt] \hline
\end{tabular}

\vspace{8pt}
\renewcommand{\arraystretch}{2.2}
\begin{tabular}{|>{\centering\arraybackslash}m{7.5cm}|>{\centering\arraybackslash}m{3.2cm}|>{\centering\arraybackslash}m{3.2cm}|}
\hline
\textbf{الاسم واللقب} & \textbf{الفوج} & \textbf{العلامة} \\ \hline
 & & \hfill /20 \\ \hline
\end{tabular}
\end{center}
\end{document}
